%!BIB program=biber

\documentclass[12pt,aspectratio=169]{beamer} %类型为文章
\usepackage[UTF8]{ctex} %中文编码宏
\usepackage{multicol} %分栏控制宏
\usepackage{hyperref} %超链接宏
\usepackage{lastpage} %总计页的宏
\usepackage{color} %颜色控制宏
\usepackage{graphicx} %图片插入宏
\usepackage{subfigure} %子图插入宏
\usepackage{animate} %动画插入宏
\usepackage{multirow} %纵向合并宏
\usepackage{makecell} %表格换行宏
\usepackage{amsmath} %公式插入宏
\usepackage{unicode-math} %公式样式宏
\usepackage{gbt7714} %国标引用宏
\usepackage{url} %网页链接宏
\usepackage{doi} %doi号宏
\usepackage{svg}
\renewcommand{\vec}[1]{\boldsymbol{#1}} %设置向量样式

\usetheme{Berlin}
\usecolortheme{beaver}

\linespread{1.2} %行距
\setlength{\parskip}{0.5em} %段落间距
\setlength{\parindent}{2em} %缩进距离

\setmathfont{Cambria Math} %设置数学公式样式
%\bibliographystyle{gbt7714-numerical} %设置参考文献样式

%\logo{\includegraphics[height=0.1\textwidth]{images/SCU_logo.pdf}}
%\setbeamertemplate{background}{\includegraphics[height=\paperheight]{images/SCU_logo.pdf}}
\setbeamertemplate{itemize items}{$\blacksquare $}
\setbeamertemplate{caption}[numbered]


\title{工作进展回报} %设置标题
\subtitle{基于Kramers-Kronig关系的相位恢复方法}
\author{Julian OU} %设置作者
\institute[SCU]{\textit{College of Physics, Sichuan University, Chengdu 610064, China}}
\date{\today} %设置日期

\begin{document}
\maketitle %插入标题

\AtBeginSection{
    \begin{frame}
        \frametitle{目录}
        \tableofcontents[currentsection,subsectionstyle=hide]
    \end{frame}
}

\AtBeginSubsection{
    \begin{frame}
        \subsectionpage
    \end{frame}
}

\section{KK关系理论}

\begin{frame}
    \begin{block}{描述}
        在x的上半平面上解析的平方可积函数$f(x)$满足方程
        \begin{align} 
            \Im[f(x)]=\frac{1}{\pi}P \int_{-\infty}^{\infty}\frac{\Re[f(x')]}{x'-x} \, \mathrm{d}x'
        \end{align}
        其中$P$为柯西主值
    \end{block}
    如果一个复变函数，已知其实部，且满足以上条件，那么就可以使用KK关系推出其虚部。
\end{frame}

\begin{frame}
    把电场写作复振动，对电场取对数，便可以把强度与相位进行分离。
    \begin{align}
        \Re[\ln E] & =\ln I /2 \\
        \Im[\ln E] & =\arg (E)
    \end{align}
    
    如果$\ln E$的实部与虚部之间满足KK关系，就可以从强度图像中恢复相位，但是$\ln E$并不总是在上半平面上解析。
\end{frame}

\section{成像实验}

\begin{frame}
    \begin{columns}
        \begin{column}{0.3\textwidth}
            \begin{figure}
                \includegraphics[height=6cm]{images/2.png}
            \end{figure}
        \end{column}
        \begin{column}{0.6\textwidth}
            \small
            \begin{table}
                \begin{tabular}[c]{rl}
                    光源            & 超辐射发光二极管             \\
                    振镜系统        & 控制入射角                   \\
                    透镜1           & 聚焦                         \\
                    condenser聚光镜 & 会聚光源，透射场照明         \\
                    objective物镜   & 显微镜物镜，消像差           \\
                    透镜2           & 产生平行光束                 \\
                    光阑            & 提供截止空间角频率           \\
                    透镜3+透镜4     & 4f系统                       \\
                    偏振片          & 保证线偏振光，基于标量场模型 \\
                    CCD             & 接收信号                     \\
                \end{tabular}
            \end{table}
        \end{column}
    \end{columns}
\end{frame}

\begin{frame}
    \begin{columns}
        \begin{column}{0.3\textwidth}
            \begin{figure}
                \includegraphics[height=6cm]{images/1.png}
            \end{figure}
        \end{column}
        \begin{column}{0.6\textwidth}
            \begin{block}{在实验上使电场的传播符合KK关系的条件}
                \begin{itemize}
                    \item 入射波的横向波矢量为光阑的截止空间角频率
                    \item 非散射光强于散射光
                \end{itemize}
            \end{block}
            
            \qquad 其次，要在实验上实现相位恢复，还需要考虑离散的强度图像是否兼容KK关系，以及数值计算方法。
        \end{column}
    \end{columns}
\end{frame}

\section{证明过程}
\begin{frame}
    考虑将电场写作：
    \begin{align}
        \varepsilon(\vec{r})=\left| E(\vec{r}) \right| \exp(-i\vec{k}_{inc}\cdot \vec{r}) \label{A}
    \end{align}
    表示具有一定倾角的入射，其中$\vec{k}_{inc}$为入射波的横向波矢，构建一个复变函数：
    \begin{align}
        \chi (\vec{r})=\ln \varepsilon(\vec{r})=\ln E(\vec{r})-i\vec{k}_{inc}\cdot \vec{r}
    \end{align}
    实部为$\ln I /2$，虚部为$\arg(E)-\vec{k}_{inc}\cdot \vec{r}$ \label{C}
\end{frame}

\begin{frame}
    考虑单色入射光是沿着$x$方向倾斜的平面波：
    \begin{align}
        E_{inc}=\left| E_{inc}(\vec{r}) \right| \exp(-i\vec{k}_{inc}\cdot \vec{r})=\left| E_{inc}(\vec{r}) \right| \exp(-i k_{inc} x)
    \end{align}
    在照射到物体以后，会产生散射光，考虑带限为$-k_0\sim k_0$的出射波：
    \begin{align}
        \varepsilon(\vec{r})=E_u(\vec{r})+E_s(\vec{r})
    \end{align}
    其中$E_u(\vec{r})=\left| E_{u}(\vec{r}) \right| \exp(-i k_{inc} x)$为与入射波具有相同频率的无散射场，$E_s$为由除了$k_{inc}$以外频率组成的散射场。
\end{frame}

\begin{frame}
    因此(\ref{A})式可以写作：
    \begin{align}
        \chi (\vec{r})=\ln \left[E_u(\vec{r})+E_s(\vec{r})\right]=\ln [1+\frac{E_s(\vec{r})}{E_u(\vec{r})}]+\ln \left| E_u \right|
    \end{align}
    其中$\ln \left| E_u \right|$为常数，不会对解析性产生影响，当非散射场强度大于散射场，即$|E_u|>|E_s|$时，对数项可以进行幂级数展开：
    \begin{align}
        \ln [1+\frac{E_s(\vec{r})}{E_u(\vec{r})}]=\sum_{n=0}^{\infty}\frac{-1^n}{n+1}\left(\frac{E_s}{E_u}\right)^{n+1} \label{B}
    \end{align}
    因此只要我们证明了$E_s/E_u$是解析的，就可以证明$\chi(\vec{r})$是解析的。
\end{frame}

\begin{frame}
    \begin{block}{Titschmarch theorem}
        如果函数$f(x)$的傅里叶变换在$k_x<0$时为零，那么$f(x)$在x的上半平面上解析。
    \end{block}

    对$E_s/E_u$作傅里叶变换（$E_u(\vec{r})=\left| E_{u}(\vec{r}) \right| \exp(-i k_{inc} x)$），得
    \begin{align}
        \mathcal{F} [\frac{E_s}{E_u}](\vec{k})=\int \frac{E_s}{|E_u|} \exp[-(\vec{k}-\vec{k}_{inc}) \cdot \vec{r}] \, \mathrm{d}\vec{r}=\frac{1}{|E_u|} \mathcal{F} [E_s](\vec{k}-k_{inc}\vec{e_{kx}})
    \end{align} 
\end{frame}

\begin{frame}
    联系到实验上，成像系统的带通受到傅里叶平面上的光阑限制，波矢空间取值范围为$-k_0\sim k_0$。
    \begin{align}
        \mathcal{F} [\frac{E_s}{E_u}](\vec{k})=\frac{1}{|E_u|} \mathcal{F} [E_s](\vec{k}-k_{inc}\vec{e_{kx}})
    \end{align} 
    如果$k_{inc}=k_0$，那么在上式中，对于$k_x<0$的区域，傅里叶变换的值就为零，便证明了$E_s/E_u$在上半平面上解析，因此证明了(\ref{B})(\ref{C})在上半平面上解析，利用KK关系，电场于是电场的表达式可以写为
    \begin{align}
        E(\vec{r})=\exp\{[\frac{\ln I(\vec{r})}{2}-\frac{i}{\pi}P\int_{-\infty}^{\infty}\frac{\ln I(\vec{r'})}{2(r'_{\parallel} -r_{\parallel})}\, \mathrm{d}\vec{r'}]+i\vec{k}_{inc} \cdot \vec{r}\}
    \end{align}
\end{frame}



\end{document}